\abstract % Структурный элемент: РЕФЕРАТ

Ключевые слова: TLA+, трассировка, валидация трасс, модельная проверка,
распределённые системы, Two-Phase Commit.

Основная цель работы — разработка и апробация подхода к валидации
реализации распределённого алгоритма на основе формальной спецификации
TLA+ с использованием анализа трасс исполнения.

В работе рассматривается задача сопоставления поведения реальной
распределённой программы со спецификацией высокого уровня. В качестве
объекта исследования выбран алгоритм двухфазной транзакции (Two-Phase Commit),
широко применяемый для обеспечения согласованности в распределённых
транзакционных системах. Реализован прототип алгоритма на языке~C,
включающий координатора (Transaction Manager), участников (Resource
Managers) и сетевой сервер сообщений.

Для связывания кода со спецификацией разработан механизм инструментирования,
позволяющий фиксировать изменения только выбранных переменных спецификации,
а не полные состояния системы. Трассы исполнения формируются в формате
NDJSON и затем автоматически объединяются и подаются на вход модели
проверки TLC. Проверка корректности сводится к задаче ограниченной
модельной проверки, где TLC анализирует соответствие полученной трассы
спецификации TLA+.

В ходе экспериментов продемонстрировано обнаружение логической ошибки в
реализации алгоритма, связанной с некорректным учётом сообщений от участников.
Использование трассировки и формальной проверки позволило надёжно выявить
расхождение между спецификацией и реализацией в условиях недетерминированной
доставки сообщений.
